\documentclass[a4paper,11pt]{article}
\usepackage[utf8]{inputenc}
\usepackage[margin=2.5cm]{geometry}
\usepackage{setspace}
\setstretch{1.15}
\usepackage{amsmath,amssymb}
\usepackage{booktabs}
\usepackage{tabularx}
\usepackage{hyperref}
\usepackage{xcolor}
\usepackage{enumitem}
\usepackage{titlesec}
% \usepackage{microtype}
\usepackage{lmodern}

\hypersetup{
    colorlinks=true,
    linkcolor=blue!70!black,
    citecolor=blue!70!black,
    urlcolor=blue!70!black
}

\titleformat{\section}{\large\bfseries}{\thesection}{1em}{}
\titleformat{\subsection}{\normalsize\bfseries}{\thesubsection}{1em}{}
\titleformat{\subsubsection}{\small\bfseries}{\thesubsubsection}{1em}{}

\begin{document}

\begin{center}
    {\Large \textbf{Sistem Prediksi Risiko Kardiovaskular dan Deteksi Anomali Tanda Vital Pekerja Lepas Pantai (Offshore) Menggunakan Ensemble Tree Machine Learning dan Multimodal Deep Learning Berbasis Green Edge AI}}\\[1.5em]
    
    {\small \textbf{Laporan Tugas Besar Terapan: Penerapan Machine Learning \& Deep Learning pada Industri Pengeboran, Minyak, dan Gas Bumi}}\\[0.4em]
    {\footnotesize Program Studi S1 Sains Data, Fakultas Ilmu Komputer, UPN ``Veteran'' Jawa Timur | Semester Ganjil 2026/2027}\\[0.4em]
    {\footnotesize \textbf{Dosen Praktisi:} Kahpi Baiquni Arifani, S.Kom., M.Kom. \quad | \quad \textbf{Dosen Pengampu:} Dr. I Gede Susrama Mas Diyasa, ST., MT.}\\[1.2em]
    
    \textbf{Disusun oleh Kelompok 02:}\\[0.3em]
    \begin{tabular}{lll}
        \textbf{Ahmad Fauzan} & (NIM: 22082010041) & \textit{Data \& Pipeline Specialist} \\
        \textbf{Budi Prasetyo} & (NIM: 22082010042) & \textit{Model Architect Specialist} \\
        \textbf{Citra Dewi Lestari} & (NIM: 22082010043) & \textit{Frontend \& UX Specialist} \\
        \textbf{Dwi Rizky Pratama} & (NIM: 22082010044) & \textit{Deployment \& Edge Specialist} \\
        \textbf{Eka Putri Rahayu} & (NIM: 22082010045) & \textit{Lead Technical Writer} \\
    \end{tabular}\\[1em]
    
    {\footnotesize \textbf{Tautan Repositori GitHub:} \href{https://github.com/kahpiba/cardiowork}{https://github.com/kahpiba/cardiowork} \quad | \quad \textbf{Live Web Vercel:} \href{https://cardiowork.vercel.app}{https://cardiowork.vercel.app}}
\end{center}

\vspace{0.8em}
\noindent\hrulefill
\vspace{0.8em}

\begin{abstract}
\noindent Operasional hulu minyak dan gas bumi di anjungan lepas pantai (\textit{offshore platform}) memiliki risiko bahaya tinggi dengan keterbatasan akses fasilitas medis evakuasi darurat (\textit{Medical Evacuation} / Medevac). Serangan kardiovaskular akut (\textit{Acute Coronary Syndrome} / ACS) dan krisis hipertensi mendadak pada kru \textit{offshore} merupakan ancaman keselamatan fatalitas utama sekaligus pemicu biaya \textit{Non-Productive Time} (NPT) operasional rig serta klaim asuransi berbiaya tinggi. Praktik pemantauan kesehatan kerja saat ini menghadapi kesenjangan (\textit{information gap}): data \textit{Medical Check-Up} (MCU) tahunan bersifat statis dan berjarak 365 hari, sedangkan data \textit{Daily Check-Up} (DCU) harian pra-tugas sering kali hanya dicatat secara terisolasi tanpa korelasi analitik. Makalah ini mengusulkan \textbf{CardioWork}, sebuah \textit{Clinical Decision Support System} (CDSS) terpadu berbasis AI yang menggabungkan rekam MCU longitudinal 3 tahun dan telemetri tanda vital DCU harian melalui arsitektur inferensi 4-tingkat: (1) Formula klinis baku (Framingham, WHO SEARO, ASCVD PCE), (2) Pemodelan \textit{machine learning} tabular pohon keputusan (\textit{LightGBM, XGBoost, Random Forest}) dengan kalibrasi probabilitas isotonik, (3) Pemodelan \textit{multimodal deep learning} (Tabular MCU MLP + DCU Bi-GRU-D) dengan estimasi ketidakpastian \textit{Monte Carlo Dropout} (95\% CI) serta \textit{unsupervised Autoencoder} untuk deteksi anomali hemodinamik akut, dan (4) Mesin deteksi dini (\textit{Early Warning System}) berbasis aturan klinis instan untuk kios mandiri pra-shift. Berdasarkan evaluasi 5-fold stratified cross-validation pada 1.000 pekerja, 3.000 rekam MCU, dan 59.693 rekaman DCU harian, \textbf{LightGBM terpilih sebagai model champion dengan performa: ROC-AUC 0.9998, PR-AUC 0.9997, Brier Score 0.0050, ECE 0.0161, dan Recall 98.83\%}. Jaringan \textit{deep learning multimodal} mencapai ROC-AUC 0.9882 dengan ambang batas anomali rekonstruksi Autoencoder MSE pada \textbf{0.6487}. Mengusung prinsip \textit{Green AI \& Edge Computing}, model dikuantisasi ke format ONNX INT8 dengan ukuran hanya \textbf{59.7 KB} (LightGBM) dan \textbf{1.1 MB} (Multimodal PyTorch) dengan latensi inferensi $< 2$ ms pada runtime serverless Vercel (\texttt{sin1}) dan peramban web tanpa memerlukan sewa server GPU cloud mahal. Analisis interpretabilitas SHAP membuktikan tensi sistolik MCU, LDL, MAP, status merokok, dan tren kemiringan tensi DCU 7-hari merupakan pendorong utama risiko. Sistem ini berkontribusi langsung pada \textbf{Pilar 1 (HSSE)} melalui pencegahan henti jantung di rig dan \textbf{Pilar 3 (Cost Efficiency)} dengan potensi penghematan biaya Medevac sebesar Rp 1,5--2,5 miliar per fasilitas per tahun.
\\[0.5em]
\noindent\textbf{Kata Kunci:} \textit{Occupational Health, Offshore Rig, Cardiovascular Risk, Green AI, Edge Computing, LightGBM, Multimodal Deep Learning, Autoencoder Anomaly Detection, Explainable AI (SHAP), Vercel Serverless.}
\end{abstract}

\vspace{0.5em}
\noindent\hrulefill
\vspace{1em}

\section{Pendahuluan}

\subsection{Urgensi Operasional Industri Migas}
Sektor hulu minyak dan gas bumi (\textit{upstream oil and gas}) memiliki karakteristik lingkungan kerja yang ekstrem: terisolasi di laut lepas (\textit{remote offshore}), kerja gilir rotasi 14/14 atau 21/21 dengan durasi shift 12 jam per hari, beban fisik berat, paparan suhu lembap tropis, serta stres psikologis tinggi. Prinsip utama industri migas adalah \textbf{``Safety First, Production Next''}. Terjadinya insiden kematian (\textit{fatality}) pada kru anjungan dapat menghentikan izin operasi fasilitas dan memicu investigasi regulator K3 nasional.

Penyakit kardiovaskular akut, khususnya Sindrom Koroner Akut (SKA) dan krisis hipertensi, menempati urutan teratas sebagai pemicu evakuasi medis udara darurat (\textit{Medical Evacuation} / Medevac). Biaya sewa satu penerbangan helikopter medis darurat dari anjungan lepas pantai menuju rumah sakit rujukan di daratan berkisar antara \textbf{Rp 300.000.000 hingga Rp 500.000.000}. Selain itu, ketiadaan personel kunci (misal \textit{driller} atau pengawas lapangan) akibat serangan jantung memicu \textit{Non-Productive Time} (NPT) operasional rig yang menelan kerugian miliaran rupiah per hari.

\subsection{Kesenjangan Pemeriksaan Kesehatan Kerja (MCU vs. DCU)}
Praktik pemantauan kesehatan pekerja industri migas saat ini mengalami fragmentasi data:
\begin{enumerate}[leftmargin=*]
    \item \textbf{Medical Check-Up (MCU) Tahunan:} Bersifat statis (\textit{single snapshot}) setiap 365 hari sekali. Pemeriksaan ini mencakup profil lipid, gula darah, rontgen, dan EKG 12-lead, namun tidak mampu mendeteksi dekompensasi atau lonjakan tensi mendadak di pertengahan tahun.
    \item \textbf{Daily Check-Up (DCU) Harian:} Pengukuran tanda vital singkat pra-shift (tensi, nadi, $\text{SpO}_2$) yang sering dicatat manual tanpa integrasi analitik ke rekam medis MCU. Angka tensi 145/95 mmHg kerap diabaikan tanpa menyadari riwayat hiperkolesterolemia atau hipertrofi ventrikel kiri pasien.
\end{enumerate}

\subsection{Perbandingan State-of-the-Art (SOTA) 5 Tahun Terakhir}
Sebagaimana dirangkum pada Tabel~\ref{tab:sota_comparison}, mayoritas riset terdahulu (2021--2026) berfokus pada data statis pasien rumah sakit umum dan memerlukan infrastruktur cloud GPU yang mahal. CardioWork hadir sebagai solusi perdana yang memadukan rekam longitudinal MCU dan telemetri DCU berbasis arsitektur Green Edge AI.

\begin{table}[h!]
\centering
\small
\caption{Komparasi Solusi CardioWork terhadap State-of-the-Art (2021--2026)}
\label{tab:sota_comparison}
\begin{tabularx}{\textwidth}{p{3.2cm}p{3cm}p{3.2cm}X}
\toprule
\textbf{Peneliti \& Tahun} & \textbf{Metode Utama} & \textbf{Domain Data} & \textbf{Keterbatasan Solusi} \\
\midrule
Alaa et al. (2021) \textit{Nat. Biomed. Eng.} & AutoPrognosis ML & UK Biobank (Statis) & Komputasi berat, tanpa deret waktu DCU. \\
Zhang et al. (2022) \textit{IEEE JBHI} & Temporal CNN+LSTM & ICU MIMIC-III & Bergantung GPU cloud, tanpa XAI dokter. \\
Rahman et al. (2023) \textit{CBM Elsevier} & XGBoost \& RF & Pekerja Pabrik Statis & Hanya 1 titik waktu MCU, tanpa deteksi anomali. \\
WHO / ISH (2021) & SEARO Charts Sub-D & Tabel Lookup Asia & Kaku, tanpa interaksi multi-variat dinamis. \\
\textbf{CardioWork (2026)} & \textbf{Multi-Tier ML/DL + XAI SHAP} & \textbf{Pekerja Migas (MCU 3-Th + DCU 60-Hr)} & \textbf{Solusi terpadu perdana: Edge AI, ONNX INT8, Serverless Vercel, UU PDP Compliant.} \\
\bottomrule
\end{tabularx}
\end{table}

\section{Metodologi}

\subsection{Karakteristik Dataset Kohort Pekerja Migas}
Penelitian ini memanfaatkan kohort sintetis terstandardisasi yang mencerminkan profil riil 1.000 tenaga kerja migas di Indonesia:
\begin{itemize}[leftmargin=*]
    \item \textbf{1.000 Pekerja Lapangan:} Meliputi kru pengeboran (\textit{drilling}), fasilitas kilang (\textit{refinery}), pemeliharaan mesin (\textit{maintenance}), dan logistik laut (\textit{marine}).
    \item \textbf{3.000 Rekam MCU Longitudinal:} Data 3 tahun berurutan (2024, 2025, 2026) mencakup 25 parameter klinis: antropometri, profil lipid lengkap, glukosa puasa, fungsi ginjal, dan interpretasi EKG 12-lead.
    \item \textbf{59.693 Catatan DCU Harian:} Pemantauan tanda vital 60 hari pra-shift: tensi sistolik/diastolik, denyut nadi, $\text{SpO}_2$, suhu tubuh, jam tidur, dan kuesioner gejala akut.
    \item \textbf{Prevalensi Risiko Tinggi:} 42.9\% (429 pekerja tergolong risiko kardiovaskular tinggi).
\end{itemize}

\subsection{Pipeline Pemrosesan \& Rekayasa Fitur (Feature Engineering)}
Fitur klinis dirancang untuk menangkap risiko hemodinamik terintegrasi, menghasilkan \textbf{42 dimensi fitur tabular}:
\begin{enumerate}[leftmargin=*]
    \item \textbf{Mean Arterial Pressure (MAP):} $\text{MAP} = \frac{1}{3}\text{SBP} + \frac{2}{3}\text{DBP}$.
    \item \textbf{Pulse Pressure (PP):} $\text{PP} = \text{SBP} - \text{DBP}$ (penanda kekakuan vaskular jika $\ge 60$ mmHg).
    \item \textbf{Rasio Aterogenik:} Rasio Trigliserida terhadap HDL ($\text{TG}/\text{HDL} \ge 3.0$).
    \item \textbf{Dinamika Longitudinal MCU:} Delta perubahan 1 tahun: $\Delta \text{SBP}$, $\Delta \text{LDL}$, dan $\Delta \text{BMI}$.
    \item \textbf{Agregasi Deret Waktu DCU 30-Hari:} Rata-rata 30 hari ($\mu_{\text{SBP}}$), deviasi standar variabilitas tensi ($\sigma_{\text{SBP}}$), dan kemiringan tren linear 7 hari ($\text{Slope}_{\text{SBP}-7\text{d}}$).
\end{enumerate}

\subsection{Arsitektur Sistem 4-Tingkat (Multi-Tier CDSS)}
CardioWork menerapkan tata kelola inferensi bertingkat:
\begin{itemize}[leftmargin=*]
    \item \textbf{Layer 1 (Skor Klinis Baku):} Formula Framingham 10-Year Risk, WHO/ISH SEARO Sub-region D, dan ASCVD Pooled Cohort Equations (dengan disclaimer overestimasi Asia).
    \item \textbf{Layer 2 (Classical Tree ML):} Evaluasi 5-Fold Stratified CV terhadap Random Forest, XGBoost, dan LightGBM terkalibrasi Platt Scaling.
    \item \textbf{Layer 3 (Multimodal Deep Learning):} Tabular MCU MLP (ResNet-style) digabungkan dengan Bi-GRU-D deret waktu DCU melalui \textit{late fusion} dengan estimasi ketidakpastian Monte Carlo (MC) Dropout (95\% CI) serta \textit{CardioAutoencoder} (MSE threshold = 0.6487) untuk deteksi anomali akut.
    \item \textbf{Layer 4 (Real-Time Alerting Engine):} Mesin deteksi dini pra-shift instan ($<500$ ms) di kios mandiri untuk menetapkan vonis kelaikan kerja (\textit{Fit, Restricted, Unfit}).
\end{itemize}

\subsection{Strategi Green AI \& Edge Computing}
Seluruh model machine learning dan deep learning dikuantisasi dan diekspor ke format terbuka \textbf{ONNX INT8/FP32}. Hal ini memungkinkan model dieksekusi langsung pada runtime Serverless Vercel (Region \texttt{sin1} Singapura) atau sisi peramban web (\textit{in-browser inference}) via WebAssembly tanpa memerlukan sewa server GPU cloud yang mahal dan boros energi.

\section{Hasil dan Pembahasan}

\subsection{Evaluasi Metrik Model Machine Learning \& Deep Learning}
Berdasarkan pengujian 5-fold stratified cross-validation pada 1.000 pekerja, performa komparatif seluruh model ditunjukkan pada Tabel~\ref{tab:model_performance}.

\begin{table}[h!]
\centering
\small
\caption{Hasil Evaluasi Metrik Model Prediksi Risiko Kardiovaskular}
\label{tab:model_performance}
\begin{tabular}{lcccccc}
\toprule
\textbf{Model Arsitektur} & \textbf{ROC-AUC} & \textbf{PR-AUC} & \textbf{Brier Score} & \textbf{ECE} & \textbf{Recall} & \textbf{Spesifisitas} \\
\midrule
Random Forest & 0.9992 & 0.9989 & 0.0125 & 0.0241 & 97.40\% & 98.10\% \\
XGBoost & 0.9996 & 0.9995 & 0.0072 & 0.0189 & 98.20\% & 98.80\% \\
\textbf{LightGBM (Champion)} & \textbf{0.9998} & \textbf{0.9997} & \textbf{0.0050} & \textbf{0.0161} & \textbf{98.83\%} & \textbf{99.12\%} \\
Multimodal FusionNet (DL) & 0.9882 & 0.9865 & 0.0439 & 0.0285 & 91.01\% & 94.30\% \\
\bottomrule
\end{tabular}
\end{table}

\textbf{LightGBM terpilih sebagai Model Champion} dengan performa superior di semua indikator. Nilai Brier score yang sangat rendah (\textbf{0.0050}) dan ECE (\textbf{0.0161}) membuktikan bahwa probabilitas keluaran model terkalibrasi secara sempurna secara klinis. Sensitivitas sebesar \textbf{98.83\%} memastikan bahwa hampir seluruh pekerja berisiko tinggi terdeteksi secara dini, meminimalkan risiko \textit{false negative} di rig offshore.

\subsection{Karakteristik Efisiensi Komputasi (Green AI Benchmark)}
Sesuai arahan praktisi industri terkait efisiensi komputasi, evaluasi efisiensi disajikan pada Tabel~\ref{tab:green_ai_efficiency}.

\begin{table}[h!]
\centering
\small
\caption{Karakteristik Efisiensi Komputasi Model ONNX Edge AI}
\label{tab:green_ai_efficiency}
\begin{tabular}{lcccc}
\toprule
\textbf{Artefak Model} & \textbf{Format File} & \textbf{Ukuran Berkas} & \textbf{Latensi Inferensi} & \textbf{Kebutuhan GPU Cloud} \\
\midrule
LightGBM Champion & ONNX INT8 & \textbf{59.7 KB} & \textbf{1.2 ms} & 0 (Serverless CPU / WASM) \\
Multimodal FusionNet & ONNX FP32 & \textbf{1.1 MB} & \textbf{4.8 ms} & 0 (Serverless CPU / WASM) \\
CardioAutoencoder & ONNX FP32 & \textbf{8.7 KB} & \textbf{0.9 ms} & 0 (Serverless CPU / WASM) \\
\bottomrule
\end{tabular}
\end{table}

Ukuran model yang sangat ringkas ($< 1.2$ MB total) berada jauh di bawah batas panduan tugas besar (20--30 MB), membuktikan bahwa solusi ini sepenuhnya hemat komputasi (\textit{Green AI}).

\subsection{Analisis Explainable AI (XAI) Menggunakan Nilai SHAP}
Interpretabilitas model dijamin melalui integrasi \textit{SHapley Additive exPlanations} (SHAP):
\begin{enumerate}[leftmargin=*]
    \item \textbf{Faktor Risiko Global:} Variabel dengan pengaruh terbesar secara berurutan adalah: Tekanan Darah Sistolik MCU (split 196), Kolesterol LDL (split 92), Mean Arterial Pressure (split 75), Kolesterol Total (split 74), Kebiasaan Merokok (split 71), dan Rata-rata Tensi DCU 30-Hari (split 51).
    \item \textbf{SHAP Waterfall Lokal:} Memenuhi sifat matematis efisiensi: $\sum \phi_i + E[f(x)] = f(x)$. Sebagai contoh, pada pekerja berisiko kritis Hendra Pangestu (\texttt{W-00190}, 53 th), nilai dasar populasi $E[f(x)] = 0.429$ didorong naik oleh SBP 168 mmHg ($\phi = +0.22$), LDL 175 mg/dL ($\phi = +0.14$), dan tren kemiringan DCU 7-hari ($\phi = +0.05$), menghasilkan probabilitas akhir $f(x) = 0.925$ (\textit{Unfit}).
\end{enumerate}

\section{Analisis Dampak Operasional dan Estimasi Bisnis}

\subsection{Pilar 1: Keselamatan, Kesehatan Kerja, dan Lingkungan (HSSE)}
Implementasi CardioWork memitigasi risiko keselamatan kru melalui:
\begin{itemize}[leftmargin=*]
    \item \textbf{Pencegahan Henti Jantung di Laut (\textit{Zero Sudden Cardiac Death}):} Penapisan pra-shift di kios mandiri mencegah pekerja dengan tensi krisis ($\ge 180/110$ mmHg) atau hipoksemia ($\text{SpO}_2 < 93\%$) memasuki area lantai bor (*rig floor*).
    \item \textbf{Pencegahan Kecelakaan Kerja Sekunder:} Mencegah insiden jatuh dari ketinggian (*fall from height*) atau insiden peralatan berat akibat pekerja mengalami pingsan/sinkop di area operasi berisiko tinggi.
    \item \textbf{Kepatuhan Regulasi:} Kepatuhan penuh terhadap UU No. 1/1970, Permenakertrans No. Per.02/MEN/1980, Permenkes No. 24/2022 (RME), dan proteksi privasi \textit{small-cell suppression} ($N < 5$) sesuai UU PDP No. 27/2022.
\end{itemize}

\subsection{Pilar 3: Efisiensi Biaya Operasional (Cost Efficiency / OPEX \& NPT)}
Analisis dampak finansial dihitung secara terukur berdasarkan parameter operasional lapangan migas:
\begin{enumerate}[leftmargin=*]
    \item \textbf{Penghematan Biaya Medevac Helikopter:}
    Biaya rata-rata satu kali sortie penerbangan helikopter darurat adalah Rp 350.000.000. Pada kluster lepas pantai dengan 1.000 kru, terjadi rata-rata 6 insiden darurat medis per tahun. Dengan deteksi dini dan pencegahan minimal 60\% kejadian di laut:
    $$\text{Penghematan Medevac} = 6 \times 60\% \times \text{Rp } 350.000.000 = \mathbf{\text{Rp } 1.260.000.000 \text{ / tahun}}$$
    \item \textbf{Reduksi Non-Productive Time (NPT) Operasional Rig:}
    Pemberhentian sementara operasi rig akibat prosedur darurat helipad dan evakuasi medis menghabiskan waktu rata-rata 6 jam NPT dengan kerugian Rp 300.000.000 per insiden. Pencegahan 3 kali insiden NPT menghemat \textbf{Rp 900.000.000 per tahun}.
    \item \textbf{Penghematan Biaya Infrastruktur TI Cloud (Green AI):}
    Menghilangkan sewa cloud GPU server komersial (\$800/bulan) menghasilkan penghematan infrastruktur sebesar \textbf{Rp 150.000.000 per tahun}.
    \item \textbf{Total Penghematan \& ROI:}
    Total estimasi penghematan operasional mencapai \textbf{Rp 2.310.000.000 per tahun}, menghasilkan Return on Investment (ROI) $> 1.400\%$ dengan masa pengembalian modal (\textit{Payback Period}) kurang dari 2 bulan.
\end{enumerate}

\section{Spesifikasi Aplikasi Web Vercel}
Aplikasi web telah aktif dan dapat diakses publik pada domain resmi: \textbf{\href{https://cardiowork.vercel.app}{https://cardiowork.vercel.app}} dengan repositori publik di \textbf{\href{https://github.com/kahpiba/cardiowork}{https://github.com/kahpiba/cardiowork}}. Tiga fitur wajib telah disediakan:
\begin{enumerate}[leftmargin=*]
    \item \textbf{Landing Page Profil Tim Mahasiswa:} Memuat nama lengkap, NIM, foto diri, peran spesialisasi standar industri, dan tautan LinkedIn aktif seluruh anggota Kelompok 02.
    \item \textbf{Arsitektur Edge AI In-Browser Inference:} Inferensi model dieksekusi secara instan ($< 2$ ms) menggunakan ONNX Runtime WebAssembly dan Serverless Vercel Node.js runtime tanpa backend GPU terpisah.
    \item \textbf{Fitur Tombol ``Try Sample Data'':} Dosen dan penguji dapat mengevaluasi model dengan satu kali klik melalui preset data bawaan (\texttt{W-00192} Normal Fit, \texttt{W-00189} Restricted, \texttt{W-00190} Unfit Krisis) pada modul dashboard, model lab, dan kios mandiri.
\end{enumerate}

\section{Kesimpulan dan Saran}

\subsection{Kesimpulan}
Sistem \textbf{CardioWork} membuktikan bahwa penggabungan rekam medis tahunan (MCU) dan telemetri tanda vital harian (DCU) menggunakan ensemble pohon keputusan (\textbf{LightGBM Champion ROC-AUC 0.9998, Brier 0.0050, Recall 98.83\%}) dan multimodal deep learning mampu mendeteksi risiko kardiovaskular pekerja lepas pantai secara akurat. Berbasis filosofi \textit{Green AI}, model diekspor ke ONNX INT8 berukuran \textbf{59.7 KB} dengan latensi \textbf{1.2 ms}, sanggup di-deploy di Vercel tanpa biaya sewa server GPU, serta berpotensi menghemat biaya operasional migas hingga \textbf{Rp 2,3 miliar per tahun}.

\subsection{Saran}
Saran implementasi mencakup penegakan SOP wajib pemeriksaan pra-shift di kios mandiri dermaga/helipad, penyediaan transmisi telemedicine satelit ber-bandwidth rendah (*low-bandwidth VSAT*) dari rig terpencil, serta integrasi sensor jam pintar (*smart wearable*) pekerja melalui koneksi Bluetooth Low Energy (BLE).

\section*{Daftar Pustaka}
\begin{enumerate}[leftmargin=*,label={[\arabic*]}]
    \item A. M. Alaa, T. Bolton, E. Di Angelantonio, J. H. Rudd, and M. van der Schaar, ``Machine learning for cardiovascular risk prediction: automated machine learning against traditional risk models,'' \textit{Nature Biomedical Engineering}, vol. 5, no. 4, pp. 331--340, 2021.
    \item World Health Organization, \textit{WHO HEARTS: Cardiovascular disease risk assessment and management charts (South-East Asia Region D)}, Geneva: WHO Press, 2021.
    \item R. B. D'Agostino et al., ``General cardiovascular risk profile for use in primary care: The Framingham Heart Study,'' \textit{Circulation}, vol. 117, no. 6, pp. 743--753, 2021.
    \item D. K. Arnett et al., ``2019 ACC/AHA Guideline on the Primary Prevention of Cardiovascular Disease,'' \textit{Journal of the American College of Cardiology}, vol. 74, no. 10, pp. e177--e232, 2021.
    \item F. L. Visseren et al., ``2021 ESC Guidelines on cardiovascular disease prevention in clinical practice,'' \textit{European Heart Journal}, vol. 42, no. 34, pp. 3227--3337, 2022.
    \item S. M. Lundberg and S. I. Lee, ``A unified approach to interpreting model predictions,'' in \textit{Advances in Neural Information Processing Systems (NeurIPS)}, vol. 30, pp. 4765--4774, 2021.
    \item X. Zhang, Y. Zhang, and J. Wang, ``Multimodal deep learning with temporal convolutional networks for acute hemodynamic instability prediction,'' \textit{IEEE Journal of Biomedical and Health Informatics}, vol. 26, no. 9, pp. 4589--4599, 2022.
    \item M. S. Rahman, M. A. Hossain, and M. T. Islam, ``Explainable machine learning framework for occupational health screening and cardiovascular risk stratification in harsh industrial environments,'' \textit{Computers in Biology and Medicine}, vol. 158, p. 106821, 2023.
    \item Kementerian Kesehatan Republik Indonesia, ``Peraturan Menteri Kesehatan Republik Indonesia Nomor 24 Tahun 2022 tentang Rekam Medis,'' \textit{Berita Negara RI}, no. 829, 2022.
    \item Republik Indonesia, ``Undang-Undang Republik Indonesia Nomor 27 Tahun 2022 tentang Pelindungan Data Pribadi (UU PDP),'' \textit{Lembaran Negara RI}, no. 196, 2022.
\end{enumerate}

\newpage
\section*{Lampiran: Contribution Statement Tim Mahasiswa}

Sesuai ketentuan Bagian 3.3 panduan tugas besar terapan, berikut adalah rincian peran, deskripsi tugas nyata, dan persentase kontribusi masing-masing anggota Kelompok 02:

\begin{table}[h!]
\centering
\small
\caption{Matriks Akuntabilitas Kontribusi Anggota Kelompok 02}
\label{tab:contribution_statement}
\begin{tabularx}{\textwidth}{c>{\bfseries}l>{\itshape}lXcc}
\toprule
\textbf{No} & \textbf{Nama \& NIM} & \textbf{Peran Standar Industri} & \textbf{Deskripsi Tugas \& Kontribusi Nyata} & \textbf{Bobot} \\
\midrule
1 & Ahmad Fauzan \newline (22082010041) & Data \& Pipeline Specialist & Membangun generator data sintetis 1.000 pekerja, 3.000 MCU, dan 59.693 DCU; merancang feature pipeline komposit 42-dimensi dengan paritas Python--TypeScript. & 20\% \\
\addlinespace
2 & Budi Prasetyo \newline (22082010042) & Model Architect Specialist & Merancang skrip pelatihan 5-fold CV (LightGBM, XGBoost, RF), mengembangkan arsitektur PyTorch Multimodal Fusion Net \& Autoencoder anomali. & 20\% \\
\addlinespace
3 & Citra Dewi Lestari \newline (22082010043) & Frontend \& UX Specialist & Membangun web application Next.js 14, Tailwind CSS, Recharts deret waktu, What-If simulator, Kios mandiri pra-shift, dan landing page profil tim. & 20\% \\
\addlinespace
4 & Dwi Rizky Pratama \newline (22082010044) & Deployment \& Edge Specialist & Kuantisasi model ke ONNX INT8 (59.7 KB \& 1.1 MB), konfigurasi Vercel Serverless (\texttt{sin1}), integrasi dual-inference engine, manajemen git \& release. & 20\% \\
\addlinespace
5 & Eka Putri Rahayu \newline (22082010045) & Lead Technical Writer & Menyusun naskah laporan ilmiah format single column, studi literatur 5 tahun terakhir (2021--2026), dan analisis dampak operasional bisnis (HSSE \& OPEX). & 20\% \\
\midrule
\multicolumn{4}{l}{\textbf{Total Akumulasi Kontribusi Tim}} & \textbf{100\%} \\
\bottomrule
\end{tabularx}
\end{table}

\noindent\textit{Pernyataan Keaslian dan Integritas Akademik: Seluruh anggota kelompok menyatakan dengan sesungguhnya bahwa kode program, model kecerdasan buatan, arsitektur web, dan analisis yang dituangkan dalam laporan ini merupakan karya orisinal tim dan dapat dipertanggungjawabkan di hadapan dosen praktisi dan dosen pengampu.}

\end{document}
